% Shared TikZ styles and the small macro set used to draw sequence diagrams.

\tikzset{
  >={Stealth[length=2.1mm, width=1.6mm]},
  every picture/.style={line cap=round, line join=round},
  % boxes
  box/.style={draw=rule, fill=white, rounded corners=1.5pt, font=\sffamily\small, align=center,
    inner xsep=5pt, inner ysep=4pt, minimum height=7mm},
  % colour variants: combine with box, or with a style built on it
  clientbox/.style={fill=panel},
  serverbox/.style={draw=serverline, fill=server},
  accentbox/.style={draw=accent, fill=accentsoft},
  layer/.style={draw=hairline, fill=white, rounded corners=3pt, inner sep=7pt},
  layerlabel/.style={font=\sffamily\scriptsize\bfseries, text=muted, anchor=north west, inner sep=3pt},
  % lines
  flow/.style={->, draw=ink, line width=0.55pt},
  softflow/.style={->, draw=muted, line width=0.5pt, dashed},
  edgelabel/.style={font=\sffamily\scriptsize, fill=white, inner sep=1.5pt, text=ink},
  % state machines
  state/.style={draw=accent, fill=accentsoft, rounded corners=7pt, font=\sffamily\small\bfseries,
    minimum height=8mm, minimum width=19mm, align=center, text=accentdark},
  endstate/.style={state, draw=rule, fill=panel, text=ink},
  startdot/.style={circle, fill=ink, inner sep=0pt, minimum size=2.2mm},
  transition/.style={->, draw=ink, line width=0.55pt},
  translabel/.style={font=\sffamily\scriptsize, fill=white, inner sep=1.5pt, align=center, text=ink},
  % flowcharts
  step/.style={box, text width=#1, inner ysep=5pt},
  step/.default=34mm,
  decision/.style={draw=accent, fill=accentsoft, diamond, aspect=2.2, inner sep=1pt,
    font=\sffamily\scriptsize, align=center, text=accentdark},
  terminal/.style={box, rounded corners=8pt, fill=panel},
  % entity-relationship diagrams
  entity/.style={rectangle split, rectangle split parts=2, draw=rule, fill=white,
    rectangle split part fill={panel, white}, font=\ttfamily\scriptsize, align=flush left,
    inner xsep=4pt, inner ysep=2.5pt, text width=#1, rounded corners=1pt},
  entity/.default=36mm,
  privateentity/.style={entity=#1, dashed},
  privateentity/.default=36mm,
  relation/.style={draw=muted, line width=0.5pt},
  card/.style={font=\sffamily\scriptsize, text=muted, inner sep=1pt},
  % sequence diagrams
  actor/.style={draw=rule, fill=panel, rounded corners=1.5pt, font=\sffamily\scriptsize\bfseries,
    align=center, minimum height=8mm, minimum width=19mm, inner xsep=3pt, anchor=south},
  serveractor/.style={actor, draw=serverline, fill=server},
  lifeline/.style={draw=hairline, line width=0.6pt},
  call/.style={->, draw=ink, line width=0.5pt},
  reply/.style={->, draw=muted, line width=0.5pt, dash pattern=on 2.2pt off 1.6pt},
  push/.style={->, draw=accent, line width=0.6pt},
  msglabel/.style={font=\sffamily\scriptsize, text=ink, above, inner sep=1.2pt, align=center, fill=white, fill opacity=0.9, text opacity=1},
  seqnote/.style={draw=hairline, fill=white, font=\sffamily\scriptsize, align=center,
    inner xsep=4pt, inner ysep=2.5pt, rounded corners=1pt, text=muted},
  fragment/.style={draw=rule, line width=0.45pt, fill=panel!45},
  fraglabel/.style={draw=rule, fill=white, font=\sffamily\scriptsize\bfseries, inner sep=2pt},
  fragguard/.style={font=\sffamily\scriptsize\itshape, text=muted, inner sep=2pt},
}

% \participant[style]{node name}{x in cm}{label}
\newcommand{\participant}[4][actor]{\node[#1] (#2) at (#3,0) {#4};}

% \lifelines{y bottom in cm}{comma-separated node names}, drawn behind everything else.
\newcommand{\lifelines}[2]{%
  \begin{scope}[on background layer]
    \foreach \p in {#2} {\draw[lifeline] (\p.south) -- (\p |- 0,-#1);}
  \end{scope}}

% \msg[style]{from}{to}{y in cm}{label}: a horizontal message between two lifelines.
\newcommand{\msg}[5][call]{\draw[#1] (#2 |- 0,-#4) -- (#3 |- 0,-#4) node[msglabel, midway] {#5};}

% \selfmsg{participant}{y in cm}{label}: a call a participant makes to itself.
\newcommand{\selfmsg}[3]{%
  \draw[call] (#1 |- 0,-#2) -- ++(0.32,0) -- ++(0,-0.3) -- ++(-0.32,0);
  \node[font=\sffamily\scriptsize, anchor=west, align=left, inner sep=1pt, fill=white, fill opacity=0.9, text opacity=1] at ($(#1 |- 0,-#2)+(0.4,-0.15)$) {#3};}

% \seqnote{left participant}{right participant}{y in cm}{text}: a note spanning lifelines.
\newcommand{\seqnote}[4]{%
  \path let \p1=(#1), \p2=(#2) in
    node[seqnote, anchor=north, minimum width={\x2-\x1+12mm}] at ({(\x1+\x2)/2}, -#3 cm) {#4};}

% \fragment{left participant}{right participant}{y top}{y bottom}{label}: a combined fragment.
\newcommand{\fragment}[5]{%
  \begin{scope}[on background layer]
    \draw[fragment] ($(#1 |- 0,-#3)+(-1.05,0)$) rectangle ($(#2 |- 0,-#4)+(1.05,0)$);
  \end{scope}
  \node[fraglabel, anchor=north west] at ($(#1 |- 0,-#3)+(-1.05,0)$) {#5};}

% \fragsep{left participant}{right participant}{y}{guard}: a dashed divider inside a fragment.
\newcommand{\fragsep}[4]{%
  \draw[rule, dash pattern=on 3pt off 2pt, line width=0.45pt]
    ($(#1 |- 0,-#3)+(-1.05,0)$) -- ($(#2 |- 0,-#3)+(1.05,0)$);
  \node[fragguard, anchor=north west] at ($(#1 |- 0,-#3)+(-1.0,0)$) {#4};}

% \fragguardat{left participant}{y}{guard}: the guard of a fragment's first operand.
\newcommand{\fragguardat}[3]{\node[fragguard, anchor=north west] at ($(#1 |- 0,-#2)+(-1.0,0)$) {#3};}
